\scanpage{22}

$f\in\Gamma(U,F)$, montrons qu'on peut prolonger $f$ en une section de $f$ sur $X$. Il suffit de prolonger dans un voisinage de $x$, grâce à l'hyp. de récurrence. Or c'est encore possible grâce à l'hyp $P_2$ en $x$, et le lem. 5. (\unsure{injectivité} du \unsure{ledit}). OK.

\uline{Définition} 6\quad Soit $\mathcal F$ une catégorie fibrée sur la catégorie des ouverts d'un espace top. $X$, soit $Y$ une partie fermée de $X$. On dit que $\mathcal F$ satisfait la condition $P_1$ \deleted{$(P_2$} (resp. $P_2$, resp $P_3$) rel. à $Y$, si pour tt ouvert $V\subset X$, le foncteur

% margin note written sideways
\margin{Si $X$ est loc. noethérien \unsure{spectral|spécial}, on dit que $\mathcal F$ satisfait la condition $P_1$ ($P_2$, $P_3$) si elle satisfait la condition $P_1$ ($P_2$, $P_3$) rel à $\{x\}$, où $x$ est le pt fermé.}

\[
\deleted{\Gamma}\mathcal F(V)\mathrel{\unsure{\to|-}}\mathcal F(V\cap U)\qquad(\text{où } U=X-Y)
\]

est fidèle (resp. pleinement fidèle, \unsure{respt|resp.} une équivalence de catégories). On dit que $\mathcal F$ satisfait la condition $P_1$ (resp. $P_2$, resp $P_3$) en un pt $x\in X$, si le foncteur

\[
\underset{W}{\underset{\longrightarrow}{\mathrm{Pseudolim}}}\ \deleted{\unsure{\mathcal F|\{}}\mathcal F(W)\mathrel{\unsure{-|\to}}
\underset{W}{\underset{\longrightarrow}{\mathrm{Pseudolim}}}\ \mathcal F(W-W\cap\overline x),
\]

(où $W$ parcourt les vois. ouverts de $x$, est \deleted{une} fidèle (resp\unsure{t.|.} pl. fidèle, resp\unsure{t.|.} une équivalence).

% archive sheet number: 21
